\section{结构分析}
\label{sec:analysis}

\paragraph{随输入变化的伪时序。}
固定栅格或十字路线按坐标排列像素；\method{}让显著性场决定闭合局部路线和低到高的树顺序。较弱区域较早进入全局历史，显著区域更接近最终带Tag读出。选择性状态更新可能保留宽泛的早期背景并强调较晚证据，但这只是归纳偏置，不是对Mamba记忆行为的保证。

\paragraph{三种关系严格分工。}
法线读取轮廓邻带，双向闭环整合一个连通分量，增广等高线树跨标量层级汇总。射线碰撞不决定父子关系，父圈摘要也不进入局部轮廓编码。树可以包含分裂和汇合，但因为它无环，同一次分裂产生的兄弟支路不会再次汇合；汇合组合的是具有独立低值来源的分支。

\paragraph{一个Sobel信号完成三个决定。}
每个均匀试探点的带符号法线Sobel响应表示哪一侧热力值升高；平均符号为整圈确定普通读取侧，平滑绝对强度决定密度，最强位置固定 $p_0$。三个决定来自同一局部测量的不同部分，规则简单，也省去了试探阶段的碰撞查询。与此同时，强变化区域既获得更多点，又靠近带Tag的闭环终点，这是一种有意的双重强调，必须与均匀采样和独立锚点进行消融比较。

\paragraph{固定预算的原图采样。}
周长决定每个圈有多少点，Sobel权重只移动这些点。累计加权弧长把固定预算集中到跨轮廓变化强的位置，同时不允许局部额外产生Token。随后直接从原图读取RGB，因此模型不会从CNN Stem自动获得局部纹理。坐标、层级、法线进度、法线长度和局部弧长支撑量告诉Mamba这些不规则像素序列的几何含义。叶子分数也按弧长加权，避免高密度区域仅因点数更多而占优。

\paragraph{代价与并行性。}
选择性扫描的计算量与实际路由样本数线性相关，包括普通和终止法线、每次轮廓扫描的两个方向、树段Token以及最终叶子Token。路线构建还包含显著性、Sobel、等高线树和最终碰撞查询；含 $V$ 个路线网格顶点的标准等高线树通常为 $O(V\log V)$ \cite{carr2003contourtrees}。大量短而不规则的序列、填充和树同步仍可能降低GPU利用率。

公平比较需要匹配Token宽度和训练方案，包含原图像素或Token数匹配的基线，并把确定性路线构建计入端到端延迟。路由Token数和渐近复杂度都不能单独代表真实运行速度。
